\Article{5}
Es wird nun leicht sein, den Multiplikatoren \(x,y,z\) etc. solche Werte zu geben, dass die Summe
\[
LLmm+L'L'm'm'+L''L''m''m''+\text{etc.}
\]
ihren kleinsten Wert annimmt. Offenbar braucht man dazu die absoluten Größen der mittleren Fehler \(m,m',m''\) etc. nicht zu kennen; ihr Verhältnis zueinander genügt. Wir führen daher an ihrer Stelle die Gewichte der Beobachtungen \(p,p',p''\) etc. ein, das heißt Zahlen, die den Quadraten \(mm,m'm',m''m''\) etc. umgekehrt proportional sind, wobei das Gewicht einer beliebig gewählten Beobachtung als Einheit angenommen wird. Die Größen \(x,y,z\) etc. sind daher so zu bestimmen, dass das Polynom in den Unbestimmten
\[
\begin{gathered}
\frac{(ax+by+cz+\text{etc.}+l)^2}{p}+\frac{(a'x+b'y+c'z+\text{etc.}+l')^2}{p'}\\[4pt]
{}+\frac{(a''x+b''y+c''z+\text{etc.}+l'')^2}{p''}+\text{etc.}
\end{gathered}
\]
seinen kleinsten Wert annimmt, was, wie wir voraussetzen wollen, bei \textit{bestimmten} Werten \(x^0,y^0,z^0\) etc. geschieht.

Führen wir die folgenden Bezeichnungen ein:
\SourceLeaf{067}
\[
\begin{aligned}
\frac{aa}{p}+\frac{a'a'}{p'}+\frac{a''a''}{p''}+\text{etc.}&=[aa]\\[4pt]
\frac{ab}{p}+\frac{a'b'}{p'}+\frac{a''b''}{p''}+\text{etc.}&=[ab]\\[4pt]
\frac{ac}{p}+\frac{a'c'}{p'}+\frac{a''c''}{p''}+\text{etc.}&=[ac]\\[4pt]
\frac{bb}{p}+\frac{b'b'}{p'}+\frac{b''b''}{p''}+\text{etc.}&=[bb]\\[4pt]
\frac{bc}{p}+\frac{b'c'}{p'}+\frac{b''c''}{p''}+\text{etc.}&=[bc]\\[4pt]
\frac{cc}{p}+\frac{c'c'}{p'}+\frac{c''c''}{p''}+\text{etc.}&=[cc]
\end{aligned}
\]
\begin{center}etc. sowie\end{center}
\[
\begin{aligned}
\frac{al}{p}+\frac{a'l'}{p'}+\frac{a''l''}{p''}+\text{etc.}&=[al]\\[4pt]
\frac{bl}{p}+\frac{b'l'}{p'}+\frac{b''l''}{p''}+\text{etc.}&=[bl]\\[4pt]
\frac{cl}{p}+\frac{c'l'}{p'}+\frac{c''l''}{p''}+\text{etc.}&=[cl]
\end{aligned}
\]
\begin{center}etc.\end{center}
so erfordert die Bedingung für ein Minimum offenbar, dass
\[
\left.\begin{aligned}
0&=[aa]x^0+[ab]y^0+[ac]z^0+\text{etc.}+[al]\\
0&=[ab]x^0+[bb]y^0+[bc]z^0+\text{etc.}+[bl]\\
0&=[ac]x^0+[bc]y^0+[cc]z^0+\text{etc.}+[cl]\\
&\text{etc.}
\end{aligned}\right\}\ (1)
\]
Nachdem hieraus die Größen \(x^0,y^0,z^0\) etc. durch Elimination abgeleitet worden sind, setzen wir
\[
\left.\begin{aligned}
ax^0+by^0+cz^0+\text{etc.}+l&=L\\
a'x^0+b'y^0+c'z^0+\text{etc.}+l'&=L'\\
a''x^0+b''y^0+c''z^0+\text{etc.}+l''&=L''\\
\text{etc.}\quad&
\end{aligned}\right\}\ (2)
\]
Danach ist diejenige Funktion der Größen \(v,v',v''\) etc. zur Bestimmung unserer Unbekannten am besten geeignet und mit der geringsten Unsicherheit behaftet, deren partielle Ableitungen in dem vorliegenden bestimmten Fall jeweils die Werte \(L,L',L''\) etc. haben; das Gewicht dieser Bestimmung, das wir mit \(P\) bezeichnen wollen, ist dann
\SourceLeaf{068}
\[
=\frac{1}{\dfrac{LL}{p}+\dfrac{L'L'}{p'}+\dfrac{L''L''}{p''}+\text{etc.}}\quad(3)
\]
oder \(\frac1P\) ist der Wert des oben angegebenen Polynoms für dasjenige System von Werten der Größen \(x,y,z\) etc., das die Gleichungen (1) erfüllt.

\Article{6}
Im vorhergehenden Artikel haben wir gezeigt, wie man diejenige Funktion \(U\) erkennt, die zur geeignetsten Bestimmung unserer Unbekannten dient; sehen wir nun, welchen \textit{Wert} die Unbekannte auf diese Weise erhält. Dieser Wert werde mit \(K\) bezeichnet; er ergibt sich also, wenn man in \(U\) die beobachteten Werte der Größen \(v,v',v''\) etc. einsetzt; durch dieselbe Einsetzung erhalte die Funktion \(u\) den Wert \(k\); schließlich sei \(\kappa\) der wahre Wert der Unbekannten, der sich somit aus den wahren Werten der Größen \(v,v',v''\) etc. ergeben würde, wenn wir diese entweder in \(U\) oder in \(u\) einsetzen könnten. Hieraus folgt also
\[
\begin{aligned}
k&=\kappa+le+l'e'+l''e''+\text{etc.}\\
K&=\kappa+Le+L'e'+L''l''+\text{etc.}
\end{aligned}
\]
und daher
\[
K=k(L-l)e+(L'-l')e'+(L''-l'')e''+\text{etc.}
\]
Setzen wir in dieser Gleichung für \(L-l,\ L'-l',\ L''-l''\) etc. die Werte aus (2) ein und setzen außerdem
\[
\left.\begin{aligned}
ae+a'e'+a''e''+\text{etc.}&=\mathfrak A\\
be+b'e'+b''e''+\text{etc.}&=\mathfrak B\\
ce+c'e'+c''e''+\text{etc.}&=\mathfrak C
\end{aligned}\right\}\ (4)
\]
etc., so haben wir
\[
K=k+\mathfrak A x^0+\mathfrak B y^0+\mathfrak C z^0\ \text{etc.}\qquad(5)
\]
Die Werte der Größen \(\mathfrak A,\mathfrak B,\mathfrak C\) etc. können wir zwar nicht nach den Formeln (4) berechnen, da die Fehler \(e,e',e''\) etc. unbekannt bleiben; es ist jedoch unmittelbar klar, dass jene nichts anderes sind als die Werte der Funktionen \(X,Y,Z\) etc., die entstehen, wenn man für \(v,v',v''\) etc. die beobachteten Werte einsetzt. Auf diese Weise liefert das System der Gleichungen (1), (3), (5) die vollständige Lösung unseres Problems, da sich unsere Bemerkung am Ende von Art. 2. darüber, dass die Berechnung der Größen \(l,l',l''\) etc. auf die beobachteten Werte der Grö\SourceLeaf{069}%
ßen \(v,v',v''\) etc. zu gründen ist, offenbar mit gleichem Recht auf die Berechnung der Größen \(a,a',a''\) etc. \(b,b',b''\) etc. etc. erstrecken lässt.

\clearpage
